\clearpage
\thispagestyle{empty}
\vspace*{\fill}
\begin{flushleft}
{\small\color{grau}
\textbf{\color{tief}Informatik von Grund auf} --- Band 22\\[2pt]
\textit{Die Kunst des Weglassens} $\cdot$ Programmieren in Go\\[14pt]

Herausgeber und Urheber: Jürgen Maier\\[2pt]
Entstanden 2026 in Zusammenarbeit mit Claude (Anthropic); Konzeption, Vorgaben, Prüfung und Herausgeberschaft: Jürgen Maier.\\[14pt]

\textbf{Lizenz}\\[2pt]
Dieses Werk ist lizenziert unter einer Creative-Commons-Lizenz \textbf{Namensnennung -- Weitergabe unter gleichen Bedingungen 4.0 International (CC BY-SA 4.0)}. Sie dürfen das Werk vervielfältigen, verbreiten, bearbeiten und auch kommerziell nutzen, sofern Sie den Urheber nennen und Bearbeitungen unter derselben Lizenz weitergeben.\\[2pt]
Lizenztext: \texttt{https://creativecommons.org/licenses/by-sa/4.0/legalcode.de}\\[14pt]

\textbf{Quellen und Fassungen}\\[2pt]
Quelltexte, Fehlermeldungen und aktuelle Fassungen aller Bände:\\
\texttt{https://github.com/juergenmaier2026/informatik-von-grund-auf}\\[2pt]
Schwesterreihe: \textit{Mathematik von Grund auf}, \texttt{https://doi.org/10.5281/zenodo.23698547}\\[14pt]

\textbf{Zitiervorschlag}\\[2pt]
Maier, J. (2026). \textit{Informatik von Grund auf, Band 22: Die Kunst des Weglassens}. Fassung 1.0.0. \texttt{https://github.com/juergenmaier2026/informatik-von-grund-auf}\\[14pt]

Satz: \LaTeX{} mit Palatino. Fassung 1.0.0, Oktober 2026. 118 Seiten.
}
\end{flushleft}
\vspace*{1.5cm}
\clearpage
